\chapter{Comparativa estructural con memorias de referencia}
\label{chap:anexo-comparativa-referencias}

\section{Criterios observados}
\label{sec:criterios-observados}

Para ajustar la memoria a una longitud y nivel de detalle más próximos a otros TFG técnicos, se revisaron varias memorias de referencia proporcionadas. La revisión se centró en estructura, densidad documental y tipos de contenido incluidos, no en copiar texto ni decisiones concretas. Los trabajos revisados presentan extensiones aproximadas entre 85 y 108 páginas, con una combinación de capítulos principales y anexos.

El patrón común observado es que la memoria no se limita a describir el resultado final. Normalmente dedica espacio a motivación, contexto tecnológico, requisitos, diseño, implementación, pruebas, planificación, costes, conclusiones y anexos. En proyectos con componente software, el diseño e implementación suelen ocupar una parte sustancial. En proyectos con despliegue o hardware, también se dedica espacio a entorno, validación y procedimientos.

La versión inicial de TransitOps concentraba mucha información en pocos capítulos. Técnicamente era coherente, pero explicaba de forma demasiado sintética decisiones que en una defensa conviene desarrollar: por qué se eligieron ciertas tecnologías, qué alternativas se descartaron, cómo se valida cada capa y cómo se opera el sistema una vez desplegado.

\section{Ajustes aplicados a TransitOps}
\label{sec:ajustes-aplicados}

A partir de esa revisión, la memoria se amplió en tres direcciones. La primera fue añadir un marco tecnológico. Este capítulo permite explicar ASP.NET Core, PostgreSQL, Docker, Terraform, AWS, GitHub Actions, CloudWatch y Secrets Manager como decisiones conectadas con los objetivos del proyecto. Así, la memoria no presupone que el lector conozca por qué cada herramienta aparece en la arquitectura.

La segunda dirección fue separar arquitectura general y diseño detallado. La arquitectura describe la vista global del sistema, mientras que el diseño detallado explica contrato HTTP, dominio, persistencia, autenticación, configuración, despliegue, validación y decisiones descartadas. Esta separación hace que el documento sea más fácil de defender: primero se presenta la forma del sistema y después se justifican las decisiones internas.

La tercera dirección fue reforzar anexos y operación. Los anexos recogen trazabilidad, evidencias, procedimientos, checklist y criterios de aceptación. Este material es especialmente adecuado para TransitOps porque el proyecto tiene una orientación DevOps fuerte. El valor no está solo en la API, sino en demostrar que puede desplegarse, observarse, recuperarse y destruirse de forma controlada.

\section{Resultado de la comparación}
\label{sec:resultado-comparacion}

Tras la ampliación, la memoria queda más cercana a las referencias en extensión y estructura. El documento mantiene el mismo alcance funcional, pero desarrolla mejor el razonamiento técnico. Esto evita una ampliación artificial basada en repetir información y, en cambio, convierte decisiones ya tomadas durante los sprints en contenido explicativo.

La comparación también confirma la utilidad de apoyar la memoria en figuras, capturas y anexos. TransitOps incorpora evidencias visuales de las validaciones cloud y mantiene el detalle operativo en anexos para que el documento pueda leerse y defenderse sin depender de explicaciones externas.

El resultado final es una memoria más completa para entrega: conserva una narrativa clara, explicita decisiones, documenta limitaciones y proporciona material operativo suficiente para una defensa técnica. La longitud aumenta sin cambiar el proyecto ni inventar funcionalidades no implementadas.
